\section{Application architecture}
\label{sec:arch}

\subsection{The system in one picture}

Learning Mocha is an Android application plus a small Express process, and the boundary between them
is not the usual client/server split. The device owns the data; the server owns a secret. Nothing in
\code{backend/} has a database, a session table or a disk write, and no part of the library reaches
it except text a user has deliberately typed into a chat message. Figure~\ref{fig:architecture} is
the whole system.

% System architecture: the device is the source of truth; the gateway only holds the key.
\begin{figure}[H]
\centering
\resizebox{\linewidth}{!}{%
\begin{tikzpicture}[
  font=\small,
  layer/.style={draw=mochaBrown, fill=mochaCream, rounded corners=3pt, align=center,
                minimum height=8mm, inner xsep=6pt, text width=#1},
  chip/.style={draw=mochaBrown!60, fill=mochaCreamDark, rounded corners=2pt, align=center,
               font=\scriptsize, minimum height=7mm, inner xsep=4pt},
  group/.style={draw=mochaBrown!35, dashed, rounded corners=4pt, inner sep=7pt},
  flow/.style={-{Stealth[length=2.2mm]}, draw=mochaDeep, thick},
  thin flow/.style={-{Stealth[length=1.8mm]}, draw=mochaBrown!70},
  node distance=5mm and 6mm,
]

% ---- device ----
\node[layer=3.0cm] (ui) {UI\\\scriptsize Fragments + XML Views};
\node[layer=3.0cm, below=of ui] (vm) {ViewModel\\\scriptsize StateFlow, ListState};
\node[layer=3.0cm, below=of vm] (repo) {Repositories\\\scriptsize Tree, Post, Search, Graph, Backup};
\node[layer=3.0cm, below=of repo] (room) {Room / SQLite\\\scriptsize 9 entities + FTS4};

\node[layer=3.4cm, right=14mm of vm, anchor=west] (engine)
     {AI Action Engine\\\scriptsize parse $\rightarrow$ validate $\rightarrow$ review $\rightarrow$ one transaction};
\node[layer=3.4cm, below=of engine, anchor=north west] (tools)
     {Context Tools\\\scriptsize 7 read-only queries};
\node[layer=3.4cm, above=of engine, anchor=south west] (chat)
     {Chat Repository\\\scriptsize sessions, SSE stream, rounds};

\node[group, fit=(ui) (room) (chat) (tools), label={[font=\scriptsize\bfseries\color{mochaDeep}]above:%
  Android device --- source of truth, fully usable offline}] (device) {};

% ---- server side ----
\node[layer=3.4cm, right=17mm of chat.east, anchor=west] (gw)
     {Express gateway\\\scriptsize holds the API key, no database};
\node[layer=3.4cm, below=14mm of gw] (llm) {DeepSeek API\\\scriptsize chat completions};

\draw[flow] (ui) -- (vm);
\draw[flow] (vm) -- (repo);
\draw[flow] (repo) -- (room);
\draw[thin flow] (vm.east) -- (engine.west);
\draw[thin flow] (chat.south) -- (engine.north);
\draw[thin flow] (engine.south) -- (tools.north);
\draw[thin flow] (tools.west) to[out=180, in=-20] (repo.east);
\draw[flow] (engine.south west) to[out=250, in=25]
      node[chip, pos=0.75, right=3pt] {approved batch} (room.east);
\draw[flow] (chat.east) -- node[chip, pos=0.5, above=1pt] {HTTPS} (gw.west);
\draw[flow] (gw) -- (llm);

\node[chip, below=3mm of room.south, anchor=north, text width=6.4cm]
     {Never leaves the device: posts, links, tags, dictionary, resources, settings, chat history};
\end{tikzpicture}%
}
\caption{Overall architecture. Knowledge lives only in Room on the device; the backend is a
stateless key-holding proxy, and every AI-originated write passes through validation and an
explicit user review before a single Room transaction applies it.}
\label{fig:architecture}
\end{figure}


The network is touched in exactly three places, all behind the AI tab: \code{GET /v1/health}, so the
chat screen can say ``AI unavailable'' before the user commits to a question; \code{POST /v1/chat};
and \code{POST /v1/chat/stream}, the same body answered as Server-Sent Events. There is no fourth.
Home, Browse, the reader, the editor, search, tags, the glossary, the graph and export/import never
open a socket, which is why offline is not a mode here --- it is the normal case, and the AI tab is
the exception that has to cope with a network. The app holds no API key
(\code{DEEPSEEK\_API\_KEY} lives only in the gateway's \code{.env}), and Android's Auto Backup is
switched off in the manifest so the OS does not upload the Room file either.

\subsection{Layers on the device}

Inside the app the structure is conventional MVVM with a repository layer, deliberately: the
interesting parts of this project are the knowledge model and the action protocol, and neither is
helped by an inventive architecture underneath. Rule~3 of \code{AGENTS.md} is the constraint the
codebase is held to --- \emph{``Layers: Fragment $\rightarrow$ ViewModel $\rightarrow$ Repository
$\rightarrow$ DAO. No DB/network calls in UI code.''}

Data flows one way. A repository exposes Room \code{Flow}s; a ViewModel combines them into exactly
one immutable UI-state data class published as a \code{StateFlow}; a Fragment collects that one flow
and renders it. Events travel back as plain method calls on the ViewModel.

\begin{lstlisting}[caption={Package layout under app/src/main/java/com/cs426/learningmocha/}]
LearningMochaApp.kt  Application: 8 lazy singletons, theme, backup reminder
MainActivity.kt      the only Activity: NavHost + bottom navigation
ui/                  14 Fragments + 1 bottom sheet, adapters, list-state views
viewmodel/           15 ViewModels, one UiState data class each
data/local/          entity/ (9 entities), dao/ (Node, Knowledge, Chat),
                     AppDatabase, KnowledgeSync, SeedData, InlineResources
data/repo/           Tree, Post, Search, Graph      data/prefs/ SettingsStore
ai/protocol/         Envelope, ActionParser, ActionValidator
ai/engine/           ActionExecutor, ContextTools, KbIndex
ai/chat/             ChatRepository
net/                 ApiClient, MochaApi, SseChatClient, SseFrames, ChatDtos
backup/              BackupRepository, BackupSnapshot, BackupReminder
util/                8 pure-Java classes (below)
\end{lstlisting}

Browse, end to end: \code{BrowseFragment} collects \code{BrowseViewModel.uiState}, which flat-maps
the current folder id over \code{TreeRepository.observeChildren(parentId)}, one indexed
\code{SELECT} in \code{NodeDao}. A drag ends in \code{ItemTouchHelper.clearView}, which calls
\code{viewModel.persistOrder(ids)} $\rightarrow$ \code{TreeRepository.reorder}, which writes the new
\code{orderIndex} values inside \code{db.withTransaction}. Room re-emits the query and the list
re-renders; the Fragment never mutates its own data.

Rules live in the repository, not the ViewModel and not the DAO. Renaming a post is the clearest
case: it refuses a title another post already holds, rewrites every \code{[[old title]]} in every
post that links to it, and re-derives those posts' link rows --- in one transaction, so an
interrupted rename cannot leave dangling links. There is one honest exception to the layering rule:
\code{TagsViewModel} reads \code{KnowledgeDao} directly, because neither \code{tags} nor
\code{post\_tags} exposes a \code{Flow} and the tag index needs a per-tag count no repository
returns. It is a commented shortcut, not a pattern.

\subsection{Why some of it is Java}

Rule~4 of \code{AGENTS.md}: Kotlin for Android-facing code, Java only for framework-free utilities in
\code{util/}. The split is chosen so the language boundary and the test boundary are the same line
--- a class that cannot import \code{android.*} is a class that can be unit-tested on the JVM in
milliseconds.

\noindent
\begin{tabularx}{\linewidth}{L{4.7cm} X}
\toprule
\textbf{Java class in \code{util/}} & \textbf{Framework-free job} \\
\midrule
\code{MarkdownLinkParser} & Extract \code{[[wiki-links]]} and YouTube URLs; rewrite a link title on
rename \\
\code{FtsQueryBuilder} & Sanitise typed input into a valid FTS4 \code{MATCH} expression, or \code{null} \\
\code{TreeRules} & \code{wouldCreateCycle} --- stops a move turning the tree into a graph \\
\code{ExportJsonWriter} \newline \code{ImportJsonReader} & Stream the \code{.mocha.json} envelope
out and read it back, with format and version checks \\
\code{TextDiff} & Line diff behind the review screen's ``what would change'' \\
\code{StreamingAnswerExtractor} & Pull readable prose from a JSON envelope still arriving \\
\code{ForceLayout} & Fruchterman--Reingold on primitive arrays, for the knowledge graph \\
\bottomrule
\end{tabularx}

\medskip

Each has a JVM unit test under \code{app/src/test/}, and the proportion works out at about 970 lines
of Java against 9{,}500 of Kotlin.

\subsection{Threading, lifecycle, and what survives what}

All asynchronous work is coroutines --- no \code{AsyncTask}, no \code{Executor}, no callbacks in the
data layer. Room's suspend and \code{Flow} DAOs already dispatch onto the database's own executor, so
\code{Dispatchers.IO} appears in only three places in the whole app (file streams in
\code{BackupRepository}, the SSE read loop in \code{SseChatClient}, the snapshot read in
\code{GraphViewModel}) and \code{Dispatchers.Default} once, for the force-directed layout. Everything
else needs no thread hop, which is itself an argument for Room.

Collection is tied to the view rather than the Fragment: every screen collects inside a
\code{repeatOnLifecycle(STARTED)} block on the view lifecycle owner, and every \code{uiState} is
published through \code{stateIn} with \code{SharingStarted.WhileSubscribed(5\_000)}, so a rotation
does not tear down and rebuild the underlying query. Views come from \code{viewBinding} held in a
nullable field cleared in \code{onDestroyView}.

That leaves two losses to handle explicitly. A \textbf{configuration change} is survived by the
ViewModel, so the current folder, an in-flight AI stream and a half-written editor draft all outlive
a rotation with nothing written to disk. \textbf{Process death} is not: seven ViewModels take a
\code{SavedStateHandle}, through which navigation arguments (\code{postId}, \code{tagId},
\code{sessionId}, \code{messageId}) arrive for free, and into which \code{BrowseViewModel}
additionally writes the folder being viewed. Anything more valuable is already in Room by then --- an
AI batch awaiting approval, for instance, lives as \code{actionsJson} on its \code{chat\_messages}
row, so it is still reviewable after a cold start.

\subsection{Navigation}

One \code{MainActivity}, one \code{NavHostFragment}, and \code{nav\_graph.xml} for the rest. Five
destinations are top-level and carry the \code{bottom\_nav\_menu} ids, which is all
\code{setupWithNavController} needs: \textbf{Home, Browse, Search, AI, Settings}. Nine more sit on
the back stack --- post reader (\code{postId}), post editor (\code{postId}, \code{parentId}),
favourites, dictionary, tag detail (\code{tagId}), knowledge graph (\code{focusPostId}, 0 = whole
library), tags, conversation (\code{sessionId}) and review changes (\code{messageId}) ---
and \code{MainActivity} hides the bottom bar for them with one
\code{addOnDestinationChangedListener} rather than a per-screen flag.

% Navigation map: one activity, five tabs, nine pushed destinations.
% The node widths are fixed in cm, so the diagram is scaled to the measure rather
% than re-laid-out whenever the body font changes.
\begin{figure}[H]
\centering
\resizebox{\linewidth}{!}{%
\begin{tikzpicture}[
  font=\scriptsize,
  tab/.style={draw=mochaBrown, fill=mochaCreamDark, rounded corners=2pt, align=center,
              minimum height=7mm, text width=1.85cm, inner sep=2pt},
  scr/.style={draw=mochaBrown!70, fill=mochaCream, rounded corners=2pt, align=center,
              minimum height=7mm, text width=2.15cm, inner sep=2pt},
  edge/.style={-{Stealth[length=1.6mm]}, draw=mochaBrown!70},
]

\node[font=\scriptsize\bfseries\color{mochaDeep}] at (-0.2,1.35) {Bottom navigation};
\node[tab] (home)     at (0.6,0.6) {Home};
\node[tab] (browse)   at (2.8,0.6) {Browse};
\node[tab] (search)   at (5.0,0.6) {Search};
\node[tab] (ai)       at (7.2,0.6) {AI};
\node[tab] (settings) at (9.4,0.6) {Settings};

\node[scr] (reader)  at (0.7,-1.4) {Post reader};
\node[scr] (graph)   at (3.2,-1.4) {Knowledge graph};
\node[scr] (favs)    at (5.7,-1.4) {Favourites};
\node[scr] (dict)    at (8.2,-1.4) {Dictionary};

\node[scr] (editor)  at (0.7,-3.0) {Post editor};
\node[scr] (tags)    at (3.2,-3.0) {Tags};
\node[scr] (tagdet)  at (5.7,-3.0) {Tag detail};
\node[scr] (chat)    at (8.2,-3.0) {Conversation};

\node[scr] (review)  at (5.7,-4.4) {Review changes};
\node[scr, draw=mochaSage] (player) at (0.7,-4.4) {YouTube sheet};

\draw[edge] (home) -- (reader);
\draw[edge] (home) -- (graph);
\draw[edge] (home) -- (favs);
\draw[edge] (home) -- (dict);
\draw[edge] (home) -- (tags);
\draw[edge] (browse) -- (reader);
\draw[edge] (search) -- (reader);
\draw[edge] (search) -- (tagdet);
\draw[edge] (ai) -- (chat);
\draw[edge] (chat) -- (review);
\draw[edge] (reader) -- (editor);
\draw[edge] (reader) -- (player);
\draw[edge] (reader) to[out=20, in=200] (graph);
\draw[edge] (tags) -- (tagdet);
\draw[edge] (tagdet) to[out=100, in=280] (reader.south east);
\end{tikzpicture}%
}
\caption{Navigation graph. One \code{MainActivity} hosts five bottom-navigation tabs; nine further
destinations sit on the back stack (the tab bar hides itself for them), and the YouTube player is a bottom
sheet rather than a destination so playback stops with the sheet. The reader, editor, favourites,
dictionary, tags, tag detail, graph, conversation and review destinations are all declared as global
actions, so any screen can open them without a dedicated edge in the graph.}
\label{fig:navmap}
\end{figure}


All nine are reached through \emph{global} actions. That matters more than it sounds: the reader is
opened from Home, Browse, Search, tag detail, the graph and the review screen, and per-screen actions
would have meant six near-identical edges and six chances to forget a transition. Only reader
$\rightarrow$ editor is local, because only the reader has an Edit button; every action shares one
180\,ms animation set from \code{res/anim/nav\_*}.

Two deliberate non-destinations: the YouTube player is a \code{BottomSheetDialogFragment}, so
dismissing the sheet destroys the \code{WebView} and stops playback with no lifecycle code of our
own, and the create/rename/move/delete dialogs are \code{MaterialAlertDialogBuilder} calls, because
none deserves a back-stack entry. Window insets are handled once, in \code{MainActivity}, which pads
the root view for system bars, cutout and IME under \code{enableEdgeToEdge()}.

\subsection{The invariant}

\begin{mochabox}[title={No AI-originated write reaches Room outside \code{ActionExecutor}}]
The model has no DAO, no repository and no database handle. It reads the library only through seven
named read-only queries in \code{ai/engine/ContextTools}, and changes it only by emitting a batch of
structured actions, which is parsed, validated against live data, shown on a review screen where each
item can be deselected, and then applied by \code{ai/engine/ActionExecutor} in a single
\code{db.withTransaction}. The codebase has exactly one call site for that executor:
\code{ReviewChangesViewModel}, on the approve button.
\end{mochabox}

The mechanism --- envelope format, validator, two-round context protocol, review screen, single-level
undo --- is \S\ref{sec:ai}. The architectural consequence is the part that belongs here: the AI is a
peer of the UI layer, not a layer beneath it, and it reaches the database through a narrower door
than the user's own screens do.
